\documentclass[11pt]{article}

% Change "review" to "final" to generate the final (sometimes called camera-ready) version.
% Change to "preprint" to generate a non-anonymous version with page numbers.
\usepackage[]{acl}
\usepackage{times}
\usepackage{latexsym}
\usepackage{dirtree}

% For proper rendering and hyphenation of words containing Latin characters (including in bib files)

% DISABLE IF WE CAN NOT USE ANOTHER FONT
% I JUST HATE THE DEFAULT FONT - Elite
\usepackage[sfdefault,lf]{carlito}
\renewcommand{\familydefault}{\sfdefault}
\usepackage[T1]{fontenc}
% For Vietnamese characters
% \usepackage[T5]{fontenc}
% See https://www.latex-project.org/help/documentation/encguide.pdf for other character sets

% This assumes your files are encoded as UTF8
\usepackage[utf8]{inputenc}
\usepackage{microtype}
\usepackage{inconsolata}
\usepackage{graphicx}

% If the title and author information does not fit in the area allocated, uncomment the following
%
%\setlength\titlebox{<dim>}
%
% and set <dim> to something 5cm or larger.

\title{Group 76 Progress Report:}


\author{Elite Lu, Ishpreet Nagi, Kenneth Ong \\
  \texttt{\{lue13,nagii,salimk4\}@mcmaster.ca} }

%\author{
%  \textbf{First Author\textsuperscript{1}},
%  \textbf{Second Author\textsuperscript{1,2}},
%  \textbf{Third T. Author\textsuperscript{1}},
%  \textbf{Fourth Author\textsuperscript{1}},
%\\
%  \textbf{Fifth Author\textsuperscript{1,2}},
%  \textbf{Sixth Author\textsuperscript{1}},
%  \textbf{Seventh Author\textsuperscript{1}},
%  \textbf{Eighth Author \textsuperscript{1,2,3,4}},
%\\
%  \textbf{Ninth Author\textsuperscript{1}},
%  \textbf{Tenth Author\textsuperscript{1}},
%  \textbf{Eleventh E. Author\textsuperscript{1,2,3,4,5}},
%  \textbf{Twelfth Author\textsuperscript{1}},
%\\
%  \textbf{Thirteenth Author\textsuperscript{3}},
%  \textbf{Fourteenth F. Author\textsuperscript{2,4}},
%  \textbf{Fifteenth Author\textsuperscript{1}},
%  \textbf{Sixteenth Author\textsuperscript{1}},
%\\
%  \textbf{Seventeenth S. Author\textsuperscript{4,5}},
%  \textbf{Eighteenth Author\textsuperscript{3,4}},
%  \textbf{Nineteenth N. Author\textsuperscript{2,5}},
%  \textbf{Twentieth Author\textsuperscript{1}}
%\\
%\\
%  \textsuperscript{1}Affiliation 1,
%  \textsuperscript{2}Affiliation 2,
%  \textsuperscript{3}Affiliation 3,
%  \textsuperscript{4}Affiliation 4,
%  \textsuperscript{5}Affiliation 5
%\\
%  \small{
%    \textbf{Correspondence:} \href{mailto:email@domain}{email@domain}
%  }
%}

\begin{document}
\maketitle
% \begin{abstract}
% \end{abstract}

% DO NOT REMOVE THIS LINE or else a lot of errors WILL APPEAR for line spacing
\hbadness=100000

\section{Introduction}

Healthcare providers deal with many vision related illnesses such as DME (Diabetic Macular Edema), CNV (Choroidal Neovascularization), and drusen that impair vision. They utilize retinal scans to detect and determine the conditions. An issue that arises is the uniqueness of each ailment, as each condition varies in subtle ways. Correctly identifying retinal diseases can prove to be a challenge in a plethora of cases and with approximately 30 million OCT scans performed each year, this process can prove time-consuming for healthcare professionals. By having a classification model that can help speed up this process, we can allow healthcare professionals to focus their attention on helping more individuals requiring their aid.

\section{Related Work}

Here, talk about the related work you encountered for your approach. Cite at least 5 references. Refer to item 2. No one has done exactly your task? Write about the most similar thing you can find. This should be around 0.25-0.5 pages.

\section{Dataset}

\subsection{Folder Structure}

When downloading the data, the dataset was partitioned as follows:

\newpage

\dirtree{%
  .1 Dataset.
  .2 test.
  .3 CNV.
  .3 DME.
  .3 DRUSEN.
  .3 NORMAL.
  .2 train.
  .3 CNV.
  .3 DME.
  .3 DRUSEN.
  .3 NORMAL.
  .2 val.
  .3 CNV.
  .3 DME.
  .3 DRUSEN.
  .3 NORMAL.
}

\subsection{Preprocessing}

Each image was loaded from the various folders as a PIL (Pillow, not the Python Imaging Library) Image and then converted to a tensor object. The pipeline varies slightly between \textit{train} when compared to both \textit{test} and \textit{val} (abbreviated to \textit{TV} pipeline for brevity). Both pipelines resized the images to 224 pixels by 224 pixels and normalized it with means [0.485, 0.456, 0.406] and standard deviation (std) of [0.229, 0.224, 0.225]. The values are the means and standard deviations for ImageNet, which is the dataset that ResNet is being trained on. The only difference is that before normalization, the \textit{train} dataset images might be flipped whereas the image data for the \textit{TV} pipeline is not. This allowed for the training set to have a more robust and diverse dataset. However, this is technically not considered an augmentation since we did not create new data (no data distribution changes).

\subsection{Annotation}
Each image is labeled with an integer to represent its corresponding disease. The images were already named as DISEASE-\#\#\#\#\#-\#.jpeg and located in their corresponding folders, which allowed for ease of labeling. The numbers for their corresponding disease is shown below:

\begin{center}
  \begin{tabular}{| c | c |}
    \hline \hline
    \textbf{CNV}    & 0 \\
    \hline
    \textbf{DME}    & 1 \\
    \hline
    \textbf{DRUSEN} & 2 \\
    \hline
    \textbf{CNV}    & 3 \\
    \hline \hline
  \end{tabular}
\end{center}

\section{Features}

\textbf{Describe any features you used for your model, or how your data was input to your model. Are you doing feature engineering or feature selection? Are you learning embeddings? Is it all part of one neural network? Refer to item 2. This may range from 0.25 pages to 0.5 pages.}


\subsection{Data Input}
Since we are doing a computer vision/image classification task, the feature were derived from the pixels of the images. The input for the model was images, which are comprised of pixels. Each pixel is denoted by three values: Red(R), Green(G), and Blue(B). As mentioned previously, each image was cropped and normalized with values corresponding the means and standard deviations for images from ImageNet.

\subsection{Learning Model}
Afterwards we used a pretrained neural network that is already trained on images and rather than training it, we are adapting it to suit the situation of identifying the diseases. ResNet-50 was used, which is trained on ImageNet. This is a type of convolutional neural network (CNN), which is a subset of neural networks that are used for image recognition and various additional visual related tasks such as processing text.

\subsection{Features}
In terms of selecting features, we did not directly do that. This is because it is not required as there no features to select. Since the pixels are the "feature", this would not be justified. As a result, feature engineering and feature selection was not needed. Same can be said for learning embeddings, as this is not a natural language processing task.

\section{Implementation}

Describe your model and implementation here. Refer to item 4. This may take around a page.

\section{Results and Evaluation}

How are you evaluating your model? What results do you have so far? What are your baselines? Refer to item 5. This may take around 0.5 pages.

\section{Feedback and Plans}

Write about your plans for the remainder of the project. This should include a discussion of the feedback you received from your TA, and how you plan to improve your approach. Reflect on your implementation and areas for improvement. Refer to item 6. This may be around 0.5 pages.

\section{Template Notes}

You can remove this section or comment it out, as it only contains instructions for how to use this template. You may use subsections in your document as you find appropriate.

\subsection{Tables and figures}

See Table~\ref{citation-guide} for an example of a table and its caption.
See Figure~\ref{fig:experiments} for an example of a figure and its caption.


\begin{figure}[t]
  \includegraphics[width=\columnwidth]{example-image-golden}
  \caption{A figure with a caption that runs for more than one line.
    Example image is usually available through the \texttt{mwe} package
    without even mentioning it in the preamble.}
  \label{fig:experiments}
\end{figure}

\begin{figure*}[t]
  \includegraphics[width=0.48\linewidth]{example-image-a} \hfill
  \includegraphics[width=0.48\linewidth]{example-image-b}
  \caption {A minimal working example to demonstrate how to place
    two images side-by-side.}
\end{figure*}


\subsection{Citations}

\begin{table*}
  \centering
  \begin{tabular}{lll}
    \hline
    \textbf{Output}           & \textbf{natbib command} & \textbf{ACL only command} \\
    \hline
    \citep{Gusfield:97}       & \verb|\citep|           &                           \\
    \citealp{Gusfield:97}     & \verb|\citealp|         &                           \\
    \citet{Gusfield:97}       & \verb|\citet|           &                           \\
    \citeyearpar{Gusfield:97} & \verb|\citeyearpar|     &                           \\
    \citeposs{Gusfield:97}    &                         & \verb|\citeposs|          \\
    \hline
  \end{tabular}
  \caption{\label{citation-guide}
    Citation commands supported by the style file.
  }
\end{table*}

Table~\ref{citation-guide} shows the syntax supported by the style files.
We encourage you to use the natbib styles.
You can use the command \verb|\citet| (cite in text) to get ``author (year)'' citations, like this citation to a paper by \citet{Gusfield:97}.
You can use the command \verb|\citep| (cite in parentheses) to get ``(author, year)'' citations \citep{Gusfield:97}.
You can use the command \verb|\citealp| (alternative cite without parentheses) to get ``author, year'' citations, which is useful for using citations within parentheses (e.g. \citealp{Gusfield:97}).

\subsection{References}

\nocite{Ando2005,andrew2007scalable,rasooli-tetrault-2015}

Many websites where you can find academic papers also allow you to export a bib file for citation or bib formatted entry. Copy this into the \texttt{custom.bib} and you will be able to cite the paper in the \LaTeX{}. You can remove the example entries.

\subsection{Equations}

An example equation is shown below:
\begin{equation}
  \label{eq:example}
  A = \pi r^2
\end{equation}

Labels for equation numbers, sections, subsections, figures and tables
are all defined with the \verb|\label{label}| command and cross references
to them are made with the \verb|\ref{label}| command.
This an example cross-reference to Equation~\ref{eq:example}. You can also write equations inline, like this: $A=\pi r^2$.


% \section*{Limitations}

\section*{Team Contributions}

Write in this section a few sentences describing the contributions of each team member. What did each member work on? Refer to item 7.

% Bibliography entries for the entire Anthology, followed by custom entries
%\bibliography{custom,anthology-overleaf-1,anthology-overleaf-2}

% Custom bibliography entries only
\bibliography{custom}

% \appendix

% \section{Example Appendix}
% \label{sec:appendix}

% This is an appendix.

\end{document}
